% SPDX-License-Identifier: Apache-2.0
%
% Ethernet in TxHDL. Built by //docs:eth. Every listing is included
% from the tree, and every printout and figure is what the build made
% by running the example.
\documentclass[journal,9pt]{IEEEtran}

\newcommand{\listingsize}{\fontsize{6}{7}\selectfont}
% SPDX-License-Identifier: Apache-2.0
%
% The house preamble, shared by every document in this directory. Loaded
% after \documentclass, before \title. Keeping it in one file is what
% stops two articles drifting apart in font, listing style or figure
% style.
% Computer Modern. IEEEtran selects Times (ptm) by default, so the face has
% to be chosen explicitly.
%
% Latin Modern is the Computer Modern variety used here, and the choice is
% forced rather than aesthetic. Under T1 encoding, plain `cmr` has no Type 1
% outlines unless cm-super is installed, and it is not in the pinned
% distribution, so pdflatex falls back to EC bitmaps and every glyph in the
% PDF becomes a Type 3 font. That was measured: the first build produced a
% document with 10 Type 3 fonts and no Type 1. Latin Modern is Computer
% Modern redrawn with a full T1 repertoire and real .pfb outlines, so the
% same design comes out scalable.
\usepackage[T1]{fontenc}
\usepackage[utf8]{inputenc}
\usepackage{lmodern}
% microtype is not in the pinned TeX distribution. emergencystretch alone
% keeps the two-column measure from producing overfull lines.
\setlength{\emergencystretch}{3em}

\usepackage{amsmath}
\usepackage{amssymb}
\usepackage{amsthm}
\usepackage{booktabs}
\usepackage{array}
% enumitem is not in the pinned TeX distribution; the list spacing below
% does what \setlist would have done.
\usepackage{float}
\usepackage{xcolor}
\usepackage{listings}
\usepackage{tikz}
\usetikzlibrary{arrows.meta,positioning,fit,backgrounds,calc}
\usepackage[hidelinks,bookmarks=true]{hyperref}

% Numbered, definition-style box for a rule the reader is meant to apply.
\theoremstyle{definition}
\newtheorem{rulebox}{Rule}
\newtheorem{example}{Example}

% Unnumbered asides.
\theoremstyle{remark}
\newtheorem*{tip}{Tip}
\newtheorem*{pitfall}{Pitfall}

% Tighter lists, in place of enumitem's \setlist.
\setlength{\topsep}{2pt}
\setlength{\itemsep}{1pt}
\setlength{\parsep}{0pt}

\newcommand{\code}[1]{\texttt{#1}}
\newcommand{\term}[1]{\emph{#1}}

% Syntax highlighting. Four roles, distinguishable in colour and still
% distinguishable in greyscale, because an IEEE article gets printed: the
% keyword is the only bold role, the comment is the only italic one, and the
% two remaining roles differ in lightness as well as in hue.
\definecolor{kwblue}{RGB}{20,60,150}
\definecolor{cmtgreen}{RGB}{90,120,90}
\definecolor{strred}{RGB}{150,50,40}
\definecolor{numgray}{gray}{0.45}
\definecolor{framegray}{gray}{0.70}
\definecolor{bgshade}{gray}{0.97}
% A darker fill for figure cells. The listing background has to stay very
% light to keep code readable; a figure cell has to be clearly shaded or the
% distinction it draws is invisible in print.
\definecolor{cellshade}{gray}{0.90}

% The listing size is the document's choice, because it depends on the
% column: a two-column article sets `\listingsize` to 6pt and keeps its
% listings within 70 columns, a one-column document keeps the body size
% and includes source files at 80. Lines are never broken; a line that does not fit
% overflows the frame, which is visible, rather than wrapping, which is
% not.
\providecommand{\listingsize}{\scriptsize}

% `'` is deliberately not a string delimiter: the language uses it for loop
% labels, as Rust does, so treating it as a quote would swallow the rest of
% the line.
\lstdefinestyle{txhdl}{
  basicstyle=\ttfamily\listingsize,
  keywordstyle=\bfseries\color{kwblue},
  commentstyle=\itshape\color{cmtgreen},
  stringstyle=\color{strred},
  numberstyle=\tiny\color{numgray},
  morekeywords={mod,use,pub,const,type,struct,enum,trait,impl,fn,async,let,mut,
    match,if,else,loop,while,for,in,break,continue,return,as,await,
    module,bus,channel,transaction,protocol,pipeline,flow,seq,config,
    port,stage,pipe,instance,connect,bind,tag,domain,blackbox,foreign,
    property,role,signal,handler,true,false,
    sequence,interface,modport,implements,package,import,var,wire,func,proc,
    case,default,next,fork,branch,join,state,fsm},
  morecomment=[l]{//},
  morestring=[b]",
  showstringspaces=false,
  columns=fullflexible,
  keepspaces=true,
  breaklines=false,
  % Framed, so a listing is bounded rather than floating in the column.
  frame=single,
  rulecolor=\color{framegray},
  framesep=3pt,
  backgroundcolor=\color{bgshade},
  xleftmargin=3pt,
  xrightmargin=3pt,
  aboveskip=6pt,
  belowskip=6pt,
  % Every listing is numbered and captioned, so it can be referred to by
  % number. `caption` without `float` numbers the listing and leaves it
  % where it was written, which is what a listing in running prose needs.
  captionpos=b,
  abovecaptionskip=2pt,
  belowcaptionskip=6pt,
}

% Figures drawn as text. Framed the same way, and with the highlighting
% switched off, because a box-and-arrow diagram is not code and half its
% words turning blue reads as an accident. Setting a style adds keys rather
% than clearing the ones already set, so the three roles are neutralised by
% hand rather than by leaving them out.
\lstdefinestyle{ascii}{
  basicstyle=\ttfamily\listingsize,
  keywordstyle=\ttfamily,
  commentstyle=\ttfamily,
  stringstyle=\ttfamily,
  columns=fullflexible,
  keepspaces=true,
  frame=single,
  rulecolor=\color{framegray},
  framesep=3pt,
  backgroundcolor=\color{bgshade},
  xleftmargin=3pt,
  xrightmargin=3pt,
  aboveskip=6pt,
  belowskip=6pt,
}
\lstset{style=txhdl}

% House TikZ styles. `shorten` lives on the arrow styles rather than on each
% arrow, so every arrow in the document gets the gap, including ones added
% later. TikZ clips a path at a node's border, so without it an arrowhead
% butts against the box with no gap at all. Keep `node distance` well above
% twice the shorten, or the arrow shrinks to nothing.
\tikzset{
  box/.style={draw,rounded corners=2pt,align=center,inner sep=4pt,
              font=\footnotesize},
  boxf/.style={box,fill=bgshade},
  lbl/.style={font=\scriptsize\itshape,align=center},
  fl/.style={-{Stealth[length=4pt]},shorten >=2.5pt,shorten <=2.5pt},
  fld/.style={fl,dashed},
}



\InputIfFileExists{buildstamp.tex}{}{}
\providecommand{\buildstamp}{Draft build (outside the Bazel flow).}

% A range of the Ethernet part, by its begin/end markers.
\newcommand{\ethpart}[3]{%
\lstinputlisting[rangebeginprefix=//\ begin\{,rangebeginsuffix=\},%
rangeendprefix=//\ end\{,rangeendsuffix=\},includerangemarker=false,%
linerange=#1-#1,caption={#2},label={#3}]{lib/parts/src/eth.rs}}

\title{Ethernet in TxHDL}

\author{Filip Filmar and Dragiša Janković%
\thanks{Both authors are independent. Filip Filmar is the
corresponding author, at f@hdlfactory.com; Dragiša Janković is
reached at linkedin.com/in/dragisa-jankovic.}%
\thanks{This is the tutorial on the Ethernet part, for a reader who
knows the AXI link and its AXI-Lite bridge~\cite{axi}. Every listing is
included from the project repository \code{HDL/txhdl} at build time,
and every printout and figure is what the build produced by running
the example. The cover~\cite{cover} lists every document.}%
\thanks{The authors used a large language model, Claude, as an
assistant in exploring the concepts and in writing and constructing
the documents and the programs; every commit of the repository says
so and carries the prompts that led to it, verbatim.}%
\thanks{\buildstamp}}

\markboth{Ethernet in TxHDL}{Filmar and Janković: Ethernet in TxHDL}

\begin{document}
\maketitle

\begin{abstract}
An Ethernet MAC turns frames into bytes on a wire and back. This one
is two units, a transmitter and a receiver, of about eighty lines each
of the lowered subset. Each speaks GMII, a byte a cycle,
and each runs on one clock, so on a board each can run on the clock of
its own side of the PHY. Both store a whole frame before acting on it,
so the channels to their client can be slow and can cross from one
clock to another. In front of them is an AXI-Lite peripheral of three
words, which a core or any AXI host reaches through the AXI-Lite
bridge. The build runs the three against a loop of wire with one frame
corrupted, and simulates the three netlists against that run under
\code{nvc} and under Verilator. For the Alinx AX7A200B there is a
hand-written echo design: the PHY's RGMII, the clocks, and a crossing
between the receive clock and the transmit clock, around the two
halves of the MAC. It goes through Vivado's synthesis, place and route
to a bitstream; it has not yet run on the board.
\end{abstract}

\section{What Is Here}
\label{sec:e-what}

The part is \code{eth} in \code{txhdl\_parts}, and it has three units.

\code{EthTx} takes a frame's bytes from a channel and sends the frame
on GMII: \code{txd}, a byte, and \code{tx\_en}, high while a frame is
on the wire. \code{EthRx} takes a frame off GMII, \code{rxd},
\code{rx\_dv} and \code{rx\_er}, and offers its bytes on a channel.
\code{EthLite} is the peripheral: it takes bytes from AXI-Lite writes
and gives them to the transmitter, and it answers AXI-Lite reads with
the bytes the receiver offers.

A byte on either channel is an \code{EthByte}, the byte and whether it
is its frame's last. A frame there is what a client means by a frame:
the destination address, the source address, the type and the
payload. The preamble, the padding and the check sequence are the
MAC's business and never reach the client.

\ethpart{byte}{A byte of a frame, between the MAC and its
client.}{lst:e-byte}

\section{The Frame on the Wire}
\label{sec:e-frame}

A frame on the wire is seven bytes of \code{0x55}, the start of frame
delimiter \code{0xd5}, the frame, and four bytes of check sequence. A
frame shorter than sixty bytes is padded with zeros to sixty before
its check sequence is worked out. Twelve idle bytes follow before the
next frame may begin.

The check sequence is a CRC-32. The register starts with every bit
set, takes each byte least significant bit first, and is complemented
at the end; the four bytes go out least significant first. A receiver
that runs the frame and its check sequence through the same register
finds \code{0xdebb20e3} in it when nothing was corrupted, so it needs
no second register to compare.

A step of the register is a shift right with the polynomial
\code{0xedb88320} added when the bit that leaves is set, and a byte is
the byte added to the register and eight such steps.
Listing~\ref{lst:e-crc} says exactly that: a function of one step, and
a function that calls it eight times.

It was not always able to say so. The lowering used to inline a
function by pasting each argument's text wherever the body read it, and one
step reads its argument twice, so the eighth call held the register's
expression $2^8 = 256$ times and the crate took minutes to compile.
The byte was written as the 32 sums the eight steps work out to
instead, each bit of the new register the exclusive or of at most
fourteen bits of the old one and of the byte, which is correct and
says nothing about what it is. A value an inlined body reads more than
once is a wire of the netlist now, issue 126, so the steps are steps
again and the netlist holds sixteen more wires. A test checks the
function against the register stepped a bit at a time in plain Rust,
on every byte from four registers.

\ethpart{crc}{A byte into the CRC-32 register: one step, and eight of
them.}{lst:e-crc}

\section{The Transmitter}
\label{sec:e-tx}

A frame on the wire cannot pause, so the transmitter stores the whole
frame before it sends any of it. While it is idle and holds no whole
frame, it takes a byte from \code{tx} whenever one is offered and puts
it in \code{frame}, a memory of 2048 bytes. The byte marked last makes
the frame \code{whole}. A whole frame starts the sender: the preamble,
the frame with the register running over it, zeros up to sixty bytes
with the register still running, the check sequence, and the gap. At
the end of the gap the sender says the frame is \code{sent}, the store
forgets it, and the next is taken.

The two are two processes. The store takes a byte a cycle, which is
a loop of one wait. The sender is the wire's protocol written as the
sequence it is: a wait for a whole frame, then a loop for the
preamble, a loop over the frame's bytes bounded by \code{fill}, a
loop for the zeros when the frame is short, the four bytes of the
check sequence, and a loop for the gap, each turn one byte and one
edge. The lowering numbers those waits into a state register of its
own and gives each loop a counter, so the \code{phase} and \code{pos}
registers the first version of this half kept by hand are gone.

The client may offer bytes as slowly as it likes, and the transmitter
holds it off while a frame is going out. That is what lets the channel
between the peripheral and the transmitter cross from one clock to
another on a board.

\ethpart{txstate}{The transmitter's state.}{lst:e-txstate}

\ethpart{tx}{The transmitter's step.}{lst:e-tx}

\section{The Receiver}
\label{sec:e-rx}

The wire cannot be held off, so the receiver stores the whole frame
too, and offers it only once its check sequence has passed. It waits
for \code{rx\_dv} with the delimiter on \code{rxd}, passing over the
preamble, then stores a byte a cycle for as long as \code{rx\_dv} stays
high, with the register running over every byte, the check sequence
included. When \code{rx\_dv} falls, the frame is good if the register
holds the residue, no byte came with \code{rx\_er}, and it is longer
than its check sequence. A good frame is offered on \code{rx} without
its check sequence, its last byte marked, as fast as the client takes
it.

Three kinds of frame are dropped and counted in \code{dropped}. A
frame whose check fails is one. A frame that arrives while the one
before it is still being offered is another, since there is nowhere to
put it. The third is a frame whose start the receiver missed: a busy
line with neither the preamble nor the delimiter on it is ignored
until \code{rx\_dv} falls.

This half too is two processes. The receiver proper watches the wire
every cycle, a loop of one wait: it hunts, stores, checks, and marks
the frame \code{full}. The offerer is the sequence: a wait for a full
frame, an edge on which the length is put on \code{rx\_len}, a loop
over the payload bounded by \code{len} less four, each turn one byte
offered and held until the client takes it, and then the frame is
\code{given} back and the receiver lets the next one in. Its state
register and its loop counter are the lowering's, as the
transmitter's are.

\ethpart{rxstate}{The receiver's state.}{lst:e-rxstate}

\ethpart{rx}{The receiver's step.}{lst:e-rx}

\section{The Peripheral}
\label{sec:e-lite}

\code{EthLite} is an AXI-Lite peripheral of three words, and it moves
one byte per transaction.

\begin{table}[htbp]
\caption{The peripheral's words.}
\label{tab:e-words}
\centering\footnotesize
\begin{tabular}{@{}lll@{}}
\toprule
Word & Access & Bits \\
\midrule
0, status & read & 0 a byte waits, 1 room to send \\
1, transmit & write & 7 to 0 the byte, 8 last \\
2, receive & read & 7 to 0 the byte, 8 last, 9 valid \\
\bottomrule
\end{tabular}
\end{table}

A write to the transmit word is answered only when the transmitter has
taken the byte, so a client that writes faster than the transmitter
takes is held off rather than losing bytes. A read of the receive word
takes the oldest byte if one waits and answers with bit 9 set; with
none waiting it answers with bit 9 clear and takes nothing. The
interrupt line is high while a byte waits.

One byte per transaction is slow for a frame of fifteen hundred bytes,
but the AXI-Lite bridge makes it cheaper than it looks. A fixed burst
to one address becomes one AXI-Lite transaction per beat at that
address, so a host sends a whole frame as one burst to the transmit
word and reads sixteen bytes as one burst from the receive word.

\ethpart{lite}{The peripheral.}{lst:e-lite}

\section{The Run}
\label{sec:e-run}

\code{ex\_eth} puts the three units behind \code{LiteBridge1}, behind
a host's tracker, with the transmitter's GMII looped into the
receiver's a cycle late. The host client sends three frames, each as
one fixed burst, and after each reads bursts of sixteen from the
receive word until the frame comes back or twelve bursts have come
back empty. The loop flips one data byte of the second frame on its
way.

The first frame is twenty bytes and comes back sixty, padded. The
second fails its check and does not come back, and the receiver counts
it. The third is seventy-two bytes and comes back as it went. The run
asserts all three outcomes, and that the transmitter sent three frames
and the receiver dropped one.

\begin{figure*}[t]
\centering
\input{eth_timing.tex}
\caption{The first frame: its bytes offered to the transmitter, the
frame whole and then sent, \code{tx\_en}, the receiver's
\code{rx\_dv}, the frame being received and then full, and the bytes
offered back, with the interrupt line.}
\label{fig:e-run}
\end{figure*}

\lstinputlisting[caption={\code{ex\_eth.rs}. Run with
\code{bazel run //lib/examples:ex\_eth}.},%
label={lst:e-ex}]{lib/examples/ex_eth.rs}

\lstinputlisting[style=ascii,caption={What \code{ex\_eth} printed when
the build ran it.},label={lst:e-out}]{out_eth.txt}

The transmitter, the receiver and the peripheral are lowered, and the
build simulates each netlist against this run under \code{nvc} and
under Verilator, at its ports. A test in the part holds the
transmitter's wire to a model of the frame, at fifty-nine, sixty and
sixty-one bytes and longer, and checks the three ways a frame is
dropped.

\section{On the Board}
\label{sec:e-board}

The Alinx AX7A200B has one gigabit port, a JL2121-N040I PHY on
RGMII. The PHY negotiates 10, 100 or 1000~Mbit and full or half duplex
by itself, from the levels on its configuration pins at power on, and
answers on its management bus at address 3. RGMII sends a GMII byte as two nibbles a cycle, the low
nibble on the clock's rising edge and the high nibble on its falling
edge, and a control line that is the valid line on the rising edge and
valid exclusive-or error on the falling edge. Two edges per cycle are
the FPGA's own primitives, \code{ODDR} out and \code{IDDR} in, which a
lowered unit does not name, so the board's part is written by hand in
\code{//eth/board}. The echo design in it sends every good frame it
receives straight back out.

\code{eth\_rgmii.v} is the RGMII wrapper. The PHY expects its transmit
clock a quarter cycle after the data changes. The wrapper does that
by clocking the data on 125~MHz and the transmit clock on the same
125~MHz a quarter cycle late, both from an MMCM; it does not rely on
the PHY's internal delay. The receive clock comes from the PHY
through a global buffer, and the receiving half of the MAC runs on
it.

\lstinputlisting[caption={\code{eth/board/eth\_rgmii.v}: RGMII to
GMII.},label={lst:e-rgmii}]{eth/board/eth_rgmii.v}

The receiver and the transmitter are on two clocks, the PHY's and the
FPGA's, so the channel between them crosses from one to the other.
\code{eth\_cdc.v} is that channel: an asynchronous FIFO of sixteen
words with the valid and ready of a lowered unit's channel port on each
side, its pointers in Gray code, each crossing to the other clock
through two flip-flops. The halves of the MAC store whole frames, so
the FIFO may hold the bytes up for a cycle or two without harm.

Each half of the MAC has a reset on the clock it runs on: the clock
manager's lock, which is also what holds the PHY in reset, through
two flip-flops in each domain, so that each release lands on an edge
of its own clock. The receive clock is the PHY's and does not run
until the PHY does, so its synchroniser holds the reset from its
initial value until the first edges arrive. A lowered module's reset
arm runs from its \code{rst} input and from nothing else (issue 509).

\lstinputlisting[caption={\code{lib/board/chan\_cdc.v}: a channel from
one clock to another.},label={lst:e-cdc}]{lib/board/chan_cdc.v}

\lstinputlisting[caption={\code{eth/board/eth\_echo.v}: the echo
design.},label={lst:e-echo}]{eth/board/eth_echo.v}

\code{bazel build //eth:echo\_synth} puts the lowered halves and the
three hand-written files through Vivado's synthesis for the board's
part, with no error and no critical warning. Its two warnings are that
the halves' frame counters are removed, since nothing on the board
reads them, and that the management pins are driven by a constant or
not at all.

\code{bazel build //eth:echo\_pnr} places and routes the design with
the board's pins and writes the bitstream. Every timing constraint is
met, with 1.491~ns of setup slack and 0.065~ns of hold slack over 6549
endpoints, and the design rule check finds nothing. The design takes
1330 LUTs, 776 of them memory for the two frames and the FIFO, 232
flip-flops and 19 pins. No input or output delay is constrained on the
RGMII pins yet, so that slack says nothing about the timing between
the FPGA and the PHY.

The constraints name the clock, the LEDs and the fifteen pins of the
PHY, from the table in the AX7A200B manual.

\lstinputlisting[caption={\code{eth/board/ax7a200\_eth.xdc}: the
board's pins.},label={lst:e-xdc}]{eth/board/ax7a200_eth.xdc}

\section{What Is Not Done}
\label{sec:e-not}

The echo design has run on the board and works. On September 19, 2026
the board linked at gigabit with a host and returned every frame sent
to it: five hundred of 1400 bytes and five hundred of four bytes, byte
for byte, the short ones padded to sixty, with the host's adapter
counting no frame with a bad check sequence.

It took two changes, both to the timing at the pins, which is issue~231.
The receive side had no delay, and the PHY adds none of its own, so the
MAC accepted nothing at all; the receive clock now goes through an MMCM
that shifts it a quarter cycle, which is 2~ns of its 8~ns period. The
transmit side had the opposite fault: the design sent its clock a
quarter cycle late, the PHY delays that clock again on its side, and
the two together put the PHY's sampling edge where the data changes.
The board received every frame and the host got most of them back
corrupted, about one bit in ten thousand, which the host's adapter
counted under \code{ethtool -S rx\_errors} while the kernel's own
counter stayed at zero. The transmit clock now leaves the FPGA aligned
with the data, which is what the vendor's own design for this board
does; its routed checkpoint, opened in Vivado, says so, and it pings
without loss over the same cable.

The I/O standard of the Ethernet pins is LVCMOS33 on both sides: the
core board's banks are at 3.3~V, and the carrier sets the PHY's RGMII
pins to 3.3~V.

\code{//eth/hil:run} is the test that produced those numbers. It runs
where the cable is, sends a frame and waits for it to come back before
sending the next, and compares what returns with what went out.

Only gigabit works. At 100 or 10~Mbit the PHY sends a nibble per
cycle, and neither half of the MAC takes a byte over two cycles. The
PHY negotiates any of the three speeds, so the port on the other end
of the cable has to be a gigabit port. The management interface is
left idle, so the PHY runs on its power-on configuration, nothing
restricts what it advertises, and nothing reads its link state.

Each half holds one frame. A receiver whose frame the client has not
yet read drops the next one, and a client that wants to keep frames
arriving back to back reads each promptly. The peripheral moves a
byte per transaction.

The peripheral is not yet on Vreteno's board, which would put it
behind the core's AXI-Lite bridge and its channels through two crossings
like the echo's.

Writing the part found three things in the lowering, each worked
around here at the time, and all three since fixed. A wire named for
a reserved word of either target
language, such as \code{next}, \code{out}, \code{begin}, \code{body},
\code{buf} or \code{byte}, produced a netlist that did not analyse;
the lowering refused such a name where it was chosen, and now escapes
it instead, writing \code{next\_rw} in both targets (issue 497). A
bit of a field of a channel's value, \code{wh.data.bit(8)}, was written
in Verilog as a bit select of a part select, which Verilog does not
allow and VHDL did not simulate as the run did; a bit of a part is now
one bit of the value, at the part's offset, so the peripheral asks for
the bit of the word it was written with. And a constant of $2^{31}$ or
more was
written as an unsized integer, which overflows Verilog's 32-bit signed
integer and lies above the natural VHDL guarantees, so the receiver
compared the complement of its register with a smaller constant. The
lowering now writes a constant at the width of what it stands beside,
sized hexadecimal in Verilog and a bit string in VHDL, and the
receiver compares the residue, \code{0xdebb\_20e3}, itself.

\begin{thebibliography}{9}

\bibitem{axi}
F.~Filmar and D.~Janković, ``AXI in TxHDL,'' project repository
\code{HDL/txhdl}, \code{//docs:axi}, 2026.

\bibitem{cover}
F.~Filmar and D.~Janković, ``TxHDL documents,'' project repository
\code{HDL/txhdl}, \code{//docs:cover}, 2026.

\end{thebibliography}

\end{document}
